\documentclass[11pt,a4paper]{article}
\usepackage[margin=1in]{geometry}
\usepackage{amsmath,amssymb,amsthm}
\usepackage{algorithm}
\usepackage{algpseudocode}
\usepackage{booktabs}
\usepackage{hyperref}

\theoremstyle{definition}
\newtheorem{definition}{\textbf{Definition}}[section]
\newtheorem{theorem}{\textbf{Theorem}}[section]
\newtheorem{lemma}[theorem]{\textbf{Lemma}}
\newtheorem{remark}{\textbf{Remark}}[section]

\title{\textbf{Differentiable Neural Memory and Neural Turing Machines}}
\author{\textbf{Team Scaibu} \\ \small Email: \texttt{jaydeep.9664920749@gmail.com} \\ \small \textbf{Architecture and Core Engine Team}}
\date{\textbf{October 2026}}

\begin{document}
\maketitle

\section{Definitions and Cognitive Mental Models}

\subsection{Formal Mathematical Definition}
\textbf{Definition (External Differentiable Memory and Soft Addressing Dynamics):}
Let $\mathcal{M}_t \in \mathbb{R}^{N \times D}$ denote an external memory matrix at time step $t$ organized into $N$ addressable memory slots, each of embedding dimension $D$. A \textbf{Differentiable Neural Memory} controller interfaces with $\mathcal{M}_t$ using continuous, end-to-end differentiable reading and writing operations.

Given a content-based query read key $\mathbf{k}_t \in \mathbb{R}^D$ and inverse temperature sharpening parameter $\beta > 0$, the controller computes normalized cosine similarity scores across memory slots:
{\boldmath
\begin{equation}
S(\mathbf{k}_t, \mathcal{M}_t[i]) = \frac{\mathbf{k}_t \cdot \mathcal{M}_t[i]}{\|\mathbf{k}_t\|_2 \|\mathcal{M}_t[i]\|_2 + \epsilon}
\end{equation}
}
The normalized content read addressing distribution $\mathbf{w}_t^r \in \Delta^{N-1}$ is given by the temperature-scaled softmax:
{\boldmath
\begin{equation}
w_t^r[i] = \frac{\exp\left( \beta S(\mathbf{k}_t, \mathcal{M}_t[i]) \right)}{\sum_{j=0}^{N-1} \exp\left( \beta S(\mathbf{k}_t, \mathcal{M}_t[j]) \right)}
\end{equation}
}
The continuous \textbf{Read Vector} $\mathbf{r}_t \in \mathbb{R}^D$ retrieved from memory is:
{\boldmath
\begin{equation}
\mathbf{r}_t = \sum_{i=0}^{N-1} w_t^r[i] \mathcal{M}_t[i]
\end{equation}
}
For writing, given write addressing weights $\mathbf{w}_t^w \in \Delta^{N-1}$, an erase vector $\mathbf{e}_t \in [0, 1]^D$, and an add vector $\mathbf{a}_t \in \mathbb{R}^D$, the memory matrix undergoes an element-wise erase and additive write:
{\boldmath
\begin{equation}
\mathcal{M}_{t+1}[i, d] = \mathcal{M}_t[i, d] \cdot \left( 1 - w_t^w[i] e_t[d] \right) + w_t^w[i] a_t[d]
\end{equation}
}

\subsection{Expanded Non-Technical Definition (Intuition for Anyone)}
\textbf{Intuitive Conceptual Definition:}
\begin{enumerate}
    \item \textbf{Giving Neural Networks Working RAM}: A regular neural network stores all knowledge inside its frozen connection weights (like read-only ROM). Differentiable Neural Memory gives the network an external scratchpad (RAM) where it can store intermediate facts, solve complex math, and retrieve memories thousands of steps later.
    \item \textbf{Soft Addressing Instead of Hard Pointers}: Traditional computers access memory by absolute address (e.g. ``Read Memory Address 0x7FFF''). But hard addresses cannot be trained by calculus (derivatives cannot flow through integer jumps). Differentiable memory uses ``soft focus'' like a flashlight: it shines light over all memory rows simultaneously.
    \item \textbf{Search by Meaning (Content-Addressable)}: You don't need to remember which row you wrote to; you simply present a search query (``Where did Alice put the keys?''), and the memory automatically highlights rows with similar meaning.
    \item \textbf{Differentiable Erasing and Writing}: Just like pencil and eraser, the network can partially erase old outdated numbers and write fresh information without losing end-to-end backpropagation gradients.
\end{enumerate}

\subsection{Child-Friendly Mental Model (Physical Metaphor)}
Imagine a giant whiteboard in a classroom:
\begin{itemize}
    \item \textbf{The Grid of Boxes}: The whiteboard has 100 magnetic slots ($N$ slots).
    \item \textbf{The Flashlight Query}: When looking for your lunchbox, you shine a soft flashlight ($\mathbf{k}$) across the wall. The box holding your lunchbox glows the brightest ($w^r$).
    \item \textbf{The Soft Magnet Read}: You scoop up a blended handful of whatever the flashlight illuminated ($\mathbf{r}$).
    \item \textbf{Erase and Draw}: When lunch is eaten, your assistant takes a damp sponge ($\mathbf{e}$) to wipe away the food picture and uses a marker ($\mathbf{a}$) to write ``Empty plate''!
\end{itemize}

\section{Core Mathematical Equations and Definitions}

\subsection{Cosine Content Addressing Metric}
{\boldmath
\begin{equation}
S(\mathbf{k}, \mathbf{m}) = \frac{\sum_{d=0}^{D-1} k_d m_d}{\sqrt{\sum_{d=0}^{D-1} k_d^2} \sqrt{\sum_{d=0}^{D-1} m_d^2} + \epsilon}
\end{equation}
}
\begin{itemize}
    \item \textbf{Formal Definition}: The bounded scale-invariant cosine similarity between the query key and the memory slot.
    \item \textbf{What It Means}: Measures angular alignment between the query intention and the stored memory vector.
    \item \textbf{Why It Is Important}: Insulates memory lookup from arbitrary magnitude drift of stored memory values.
\end{itemize}

\subsection{Sharpened Softmax Memory Attention Distribution}
{\boldmath
\begin{equation}
w^r[i] = \frac{\exp\left( \beta S(\mathbf{k}, \mathcal{M}[i]) - \max_j \beta S(\mathbf{k}, \mathcal{M}[j]) \right)}{\sum_{l=0}^{N-1} \exp\left( \beta S(\mathbf{k}, \mathcal{M}[l]) - \max_j \beta S(\mathbf{k}, \mathcal{M}[j]) \right)}
\end{equation}
}
\begin{itemize}
    \item \textbf{Formal Definition}: The numerically stable softmax distribution over memory slot coordinates.
    \item \textbf{What It Means}: Assigns normalized attention weights summing to 1 across all memory rows.
    \item \textbf{Why It Is Important}: Temperature $\beta > 0$ controls the trade-off between diffuse memory blending ($\beta \to 0$) and sharp single-slot focus ($\beta \to \infty$).
\end{itemize}

\subsection{Soft Differentiable Read Operator}
{\boldmath
\begin{equation}
\mathbf{r} = \sum_{i=0}^{N-1} w^r[i] \mathcal{M}[i] = (\mathbf{w}^r)^T \mathcal{M}
\end{equation}
}
\begin{itemize}
    \item \textbf{Formal Definition}: The convex combination of memory vectors weighted by addressing probabilities.
    \item \textbf{What It Means}: Reads an interpolated representation of stored memories with unbroken gradient flow.
    \item \textbf{Why It Is Important}: Enables the controller to read memory continuously without discrete non-differentiable branching.
\end{itemize}

\subsection{Composite Erase and Additive Write Operator}
{\boldmath
\begin{equation}
\mathcal{M}_{\text{next}}[i, d] = \mathcal{M}[i, d] (1 - w^w[i] e[d]) + w^w[i] a[d]
\end{equation}
}
\begin{itemize}
    \item \textbf{Formal Definition}: The bounded element-wise state mutation operator parameterized by erase vector $\mathbf{e} \in [0, 1]^D$ and add vector $\mathbf{a} \in \mathbb{R}^D$.
    \item \textbf{What It Means}: Selectively attenuates stale values in addressed slots before injecting fresh additive content.
    \item \textbf{Why It Is Important}: Provides a stable gating mechanism that prevents memory unbounded accumulation and catastrophic saturation.
\end{itemize}

\section{Rigorous Mathematical Analysis Checklist}

\begin{itemize}
    \item \textbf{Notation}: $\mathcal{M}$ (memory tensor), $\mathbf{k}$ (read key), $\beta$ (sharpness scale), $\mathbf{w}^r$ (read attention weights), $\mathbf{w}^w$ (write attention weights), $\mathbf{e}$ (erase vector), $\mathbf{a}$ (add vector).
    \item \textbf{Epistemic Status}: Established by Graves, Wayne, and Danihelka (2014) in Neural Turing Machines; extended in Differentiable Neural Computers (DNC, 2016).
    \item \textbf{Dependency Graph}: $(\mathbf{k}, \mathcal{M}) \to \mathbf{s} \to \mathbf{w}^r \to \mathbf{r}$; $(\mathbf{w}^w, \mathbf{e}, \mathbf{a}, \mathcal{M}) \to \mathcal{M}_{\text{next}}$.
    \item \textbf{Dimensions}: Memory $\mathcal{M} \in \mathbb{R}^{N \times D}$, keys $\mathbf{k} \in \mathbb{R}^D$, attention $\mathbf{w} \in \Delta^{N-1}$, erase $\mathbf{e} \in [0, 1]^D$, add $\mathbf{a} \in \mathbb{R}^D$.
    \item \textbf{Regimes}: Algorithmic reasoning, program synthesis, long-term question answering, graph traversal tasks.
    \item \textbf{Invariants}: Address probability simplex $\sum_{i=0}^{N-1} w_i = 1$, $w_i \ge 0$.
    \item \textbf{Isomorphisms}: Attention in Transformers can be viewed as an unrolled, read-only variant of differentiable neural memory where keys and values correspond to input sequence tokens.
\end{itemize}

\section{Theoretical Theorems and Formal Proofs}

\begin{theorem}[Convex Hull Containment of the Read Vector]
Let $\mathcal{M} \in \mathbb{R}^{N \times D}$ denote the memory matrix with slot vectors $\mathbf{m}_0, \dots, \mathbf{m}_{N-1}$. For any query key $\mathbf{k}$ and sharpness parameter $\beta > 0$, the synthesized read vector $\mathbf{r}$ lies strictly within the convex hull of the memory slots:
{\boldmath
\begin{equation}
\mathbf{r} \in \operatorname{conv}(\{\mathbf{m}_0, \mathbf{m}_1, \dots, \mathbf{m}_{N-1}\})
\end{equation}
}
Consequently, $\|\mathbf{r}\|_2 \le \max_{i} \|\mathbf{m}_i\|_2$.
\end{theorem}

\begin{proof}
By definition of the softmax addressing distribution:
{\boldmath
\begin{equation}
w^r[i] = \frac{\exp(\beta S_i)}{\sum_{j=0}^{N-1} \exp(\beta S_j)}
\end{equation}
}
Since $\exp(x) > 0$ for all real $x$, every weight satisfies $w^r[i] > 0$.
Furthermore:
{\boldmath
\begin{equation}
\sum_{i=0}^{N-1} w^r[i] = \frac{\sum_{i=0}^{N-1} \exp(\beta S_i)}{\sum_{j=0}^{N-1} \exp(\beta S_j)} = 1
\end{equation}
}
Thus $\mathbf{w}^r$ belongs to the interior of the probability simplex $\Delta^{N-1}$.
The read vector is defined as the linear combination:
{\boldmath
\begin{equation}
\mathbf{r} = \sum_{i=0}^{N-1} w^r[i] \mathbf{m}_i
\end{equation}
}
By definition of a convex hull, any point expressed as $\sum_i \alpha_i \mathbf{x}_i$ with $\alpha_i \ge 0$ and $\sum_i \alpha_i = 1$ is an element of $\operatorname{conv}(\{\mathbf{x}_i\})$. Therefore $\mathbf{r} \in \operatorname{conv}(\{\mathbf{m}_0, \dots, \mathbf{m}_{N-1}\})$.
By convexity of norms:
{\boldmath
\begin{equation}
\|\mathbf{r}\|_2 = \left\| \sum_{i=0}^{N-1} w^r[i] \mathbf{m}_i \right\|_2 \le \sum_{i=0}^{N-1} w^r[i] \|\mathbf{m}_i\|_2 \le \max_{0 \le i \le N-1} \|\mathbf{m}_i\|_2 \sum_{i=0}^{N-1} w^r[i] = \max_i \|\mathbf{m}_i\|_2
\end{equation}
}
This completes the proof.
\end{proof}

\begin{theorem}[Non-Negativity Preservation under Bounded Erase and Add]
Let $\mathcal{M}_0 \in \mathbb{R}_{\ge 0}^{N \times D}$ be non-negative. If the erase vector satisfies $\mathbf{e}_t \in [0, 1]^D$, the add vector satisfies $\mathbf{a}_t \in \mathbb{R}_{\ge 0}^D$, and write weights satisfy $\mathbf{w}_t^w \in [0, 1]^N$, then the updated memory matrix remains non-negative for all $t \ge 0$:
{\boldmath
\begin{equation}
\mathcal{M}_{t+1}[i, d] \ge 0, \quad \forall i \in \{0, \dots, N-1\}, \; d \in \{0, \dots, D-1\}
\end{equation}
}
\end{theorem}

\begin{proof}
Consider the update equation for slot $i$ and dimension $d$:
{\boldmath
\begin{equation}
\mathcal{M}_{t+1}[i, d] = \mathcal{M}_t[i, d] \left( 1 - w_t^w[i] e_t[d] \right) + w_t^w[i] a_t[d]
\end{equation}
}
Since $w_t^w[i] \in [0, 1]$ and $e_t[d] \in [0, 1]$, the product satisfies $0 \le w_t^w[i] e_t[d] \le 1$.
Consequently:
{\boldmath
\begin{equation}
1 - w_t^w[i] e_t[d] \ge 0
\end{equation}
}
By induction hypothesis, $\mathcal{M}_t[i, d] \ge 0$, which implies:
{\boldmath
\begin{equation}
\mathcal{M}_t[i, d] \left( 1 - w_t^w[i] e_t[d] \right) \ge 0
\end{equation}
}
Furthermore, since $w_t^w[i] \ge 0$ and $a_t[d] \ge 0$, the additive contribution is non-negative:
{\boldmath
\begin{equation}
w_t^w[i] a_t[d] \ge 0
\end{equation}
}
The sum of two non-negative terms is non-negative, hence $\mathcal{M}_{t+1}[i, d] \ge 0$.
\end{proof}

\section{Foundational Architectural Questions and Epistemic Meta-Questions}

\subsection{Architectural Question 1: Soft vs Hard Pointer Addressing}
\textbf{Question:} Why not use hard discrete memory reads (e.g. $\arg\max_i S_i$) to prevent memory slot blurring?\\
\textbf{Answer:} Discrete $\arg\max$ has zero gradients almost everywhere, preventing backpropagation from training the controller to produce proper query keys. Soft addressing allows end-to-end gradient flow. As training progresses, the sharpness parameter $\beta$ can be annealed upwards, naturally driving the soft distribution toward discrete one-hot addressing.

\textbf{Meta-Question:} How does the memory handle address collisions when multiple slots store similar information?\\
\textbf{Answer:} Content addressing alone can suffer from collisions. Advanced architectures (such as DNC) combine content addressing with temporal linkage matrices (tracking which slot was written after which) and dynamic allocation vectors (tracking unused slots).

\subsection{Architectural Question 2: Memory Footprint vs Sequence Length}
\textbf{Question:} How does differentiable neural memory compare to Transformer self-attention across long sequences?\\
\textbf{Answer:} Standard Transformer self-attention caches all historical key-value pairs, causing quadratic $\mathcal{O}(T^2)$ computational complexity and unbounded memory growth as sequence length $T \to \infty$. Differentiable neural memory maintains a fixed-size memory buffer $\mathcal{M} \in \mathbb{R}^{N \times D}$, maintaining constant $\mathcal{O}(N D)$ memory footprint regardless of whether $T = 100$ or $T = 1,000,000$.

\textbf{Meta-Question:} Can differentiable memory learn complex algorithmic subroutines like sorting or pathfinding?\\
\textbf{Answer:} Yes. In Graves et al. (2016), DNCs learned to navigate shortest paths in graph structures (such as the London Underground) and execute associative recall tasks from scratch using purely gradient-based optimization.

\section{Algorithmic Implementation Specification}

\begin{algorithm}[H]
\caption{Differentiable Neural Memory Content Read and Gated Write}
\begin{algorithmic}[1]
\Require Memory matrix $\mathcal{M} \in \mathbb{R}^{N \times D}$, read key $\mathbf{k} \in \mathbb{R}^D$, key sharpness $\beta > 0$, write weights $\mathbf{w}^w \in \mathbb{R}^N$ (optional), erase vector $\mathbf{e} \in [0, 1]^D$ (optional), add vector $\mathbf{a} \in \mathbb{R}^D$ (optional)
\Ensure Read vector $\mathbf{r} \in \mathbb{R}^D$, read weights $\mathbf{w}^r \in \mathbb{R}^N$, updated memory $\mathcal{M}_{\text{next}} \in \mathbb{R}^{N \times D}$, memory norm $\|\mathcal{M}_{\text{next}}\|_F$
\State $\epsilon \gets 10^{-8}$
\State $norm\_k \gets \sqrt{\sum_{d=0}^{D-1} \mathbf{k}[d]^2} + \epsilon$
\State Initialize array $\mathbf{s} \in \mathbb{R}^N$
\For{$i = 0$ \textbf{to} $N-1$} \Comment{Compute Cosine Similarities}
    \State $dot \gets \sum_{d=0}^{D-1} \mathbf{k}[d] \cdot \mathcal{M}[i][d]$
    \State $norm\_m \gets \sqrt{\sum_{d=0}^{D-1} \mathcal{M}[i][d]^2} + \epsilon$
    \State $\mathbf{s}[i] \gets \beta \cdot (dot / (norm\_k \cdot norm\_m))$
\EndFor
\State $s_{\max} \gets \max_i \mathbf{s}[i]$
\State $sum\_exp \gets \sum_{i=0}^{N-1} \exp(\mathbf{s}[i] - s_{\max})$
\State Initialize $\mathbf{w}^r \in \mathbb{R}^N, \quad \mathbf{r} \in \mathbb{R}^D$
\For{$i = 0$ \textbf{to} $N-1$}
    \State $\mathbf{w}^r[i] \gets \exp(\mathbf{s}[i] - s_{\max}) / sum\_exp$
\EndFor
\For{$d = 0$ \textbf{to} $D-1$} \Comment{Perform Soft Read}
    \State $\mathbf{r}[d] \gets \sum_{i=0}^{N-1} \mathbf{w}^r[i] \cdot \mathcal{M}[i][d]$
\EndFor
\State Initialize $\mathcal{M}_{\text{next}} \in \mathbb{R}^{N \times D}, \quad norm\_sq \gets 0.0$
\For{$i = 0$ \textbf{to} $N-1$} \Comment{Optional Gated Erase and Write}
    \For{$d = 0$ \textbf{to} $D-1$}
        \State $v \gets \mathcal{M}[i][d]$
        \If{write operations enabled}
            \State $e_d \gets \max(0.0, \min(1.0, \mathbf{e}[d]))$
            \State $v \gets v \cdot (1.0 - \mathbf{w}^w[i] \cdot e_d) + \mathbf{w}^w[i] \cdot \mathbf{a}[d]$
        \EndIf
        \State $\mathcal{M}_{\text{next}}[i][d] \gets v$
        \State $norm\_sq \gets norm\_sq + v^2$
    \EndFor
\EndFor
\State \Return $(\mathbf{r}, \mathbf{w}^r, \mathcal{M}_{\text{next}}, \sqrt{norm\_sq})$
\end{algorithmic}
\end{algorithm}

\section{Big-$\Theta$ Complexity Analysis}

\begin{itemize}
    \item \textbf{Content Addressing Time Complexity}: $\Theta(N \cdot D)$ scalar multiplications to compute cosine similarities across all $N$ slots.
    \item \textbf{Soft Read Aggregation Complexity}: $\Theta(N \cdot D)$ operations to evaluate the weighted vector sum.
    \item \textbf{Memory Update Complexity}: $\Theta(N \cdot D)$ operations to execute element-wise erase and additive write.
    \item \textbf{Space Complexity}: $\Theta(N \cdot D)$ for memory buffers and intermediate addressing tensors.
    \item \textbf{Parallel Depth}: $\mathcal{D} = \Theta(\log N + \log D)$ on massively parallel accelerators.
\end{itemize}

\section{Single Canonical Repository Reference}

The complete deterministic scalar implementation, verification suite, and unit test benchmarks are published at:
\begin{center}
\url{https://github.com/Chief-Strategist-J/llm-obs-infra}
\end{center}

\end{document}
